\section{Threat Model}\label{sec:threat}
\paragraph{Assets and adversary}
The private target is the user's requested task, here one of six support actions. The defender controls the Agent and its controller. Tool providers are independent actors; the evaluated provider explicitly supplies the capabilities in Table~\ref{tab:service-contract}, without granting general service control. We consider three observers, distinguished by their rights. A \emph{transit observer} sees connection destinations, directions, ciphertext lengths and times on the links it watches. A \emph{recipient} additionally reads the requests addressed to it, including arguments. An \emph{active auditor} also chooses probes. Identifying a requested action does not establish whether it was authorized; that needs the principal and policy context. Table~\ref{tab:deployment-access} maps these rights to common deployments and to the existing evidence for each view.

\begin{table*}[t]\centering\small
\caption{Deployment and access map. The configurations are analyst-defined; cited settings test the listed view, not every component of the deployment. A gateway's plaintext access requires TLS termination.}\label{tab:deployment-access}
\setlength{\tabcolsep}{4pt}
\begin{tabularx}{\textwidth}{@{}P{.18\textwidth}P{.25\textwidth}P{.28\textwidth}Y@{}}\toprule
Deployment & Observer and permitted view & Existing evidence for that view & Protection consequence \\\midrule
Enterprise egress & A transit monitor sees routed model/tool flows, sizes, times and any exposed destination association. & WRA: supplied research-domain traces; ANT and cloud benchmark: model/combined traffic~\cite{jeong2026wra,li2026ant,pouryousef2026traffic}. & Specify which services cross the monitored egress. A union of flows is available only when the monitor can observe them. \\\addlinespace[3pt]
Shared gateway or relay & Transit sees multiplexed gateway output. A terminating gateway can additionally read addressed requests. & Live OHTTP: passive-link inference under stated relay trust~\cite{dangol2026workflow}; PrivacyTrace: semantic recipient/gateway access~\cite{duan2026privacytrace}. & Separate relay association protection from content trust. Connection sharing changes the observation unit. \\\addlinespace[3pt]
Local model, remote tools & An external observer sees remote tool traffic; the local model call has no corresponding external network leg. & WRA's destination-recovery mechanism is relevant to tool egress~\cite{jeong2026wra}; whole-deployment protection is untested in the mapped comparison. & Evaluate tool and downstream calls. Do not import a remote-model attack score into an absent network leg. \\\addlinespace[3pt]
Cloud model, remote tools & An authorized vantage may see either or both encrypted legs. A contacted tool sees its own arguments. & ANT/cloud benchmark test model traffic~\cite{li2026ant,pouryousef2026traffic}; PrivacyTrace tests recipient-visible representations~\cite{duan2026privacytrace}. & Check permitted model and tool views after the intervention; recipients and passive monitors need different protections. \\\bottomrule
\end{tabularx}
\end{table*}


\paragraph{Where cover can terminate}
Transport cover ends at an endpoint the defender controls, after the observer and before the application~\cite{apthorpe2019stp}. A relay outside a monitored egress can serve as that endpoint, but the hop from the relay to a third-party service remains. On that hop, the service's own protocol stack still offers two terminators the defender can use without the provider's cooperation. TLS~1.3 lets the sender pad any record and send padding-only records, which the receiver's TLS layer discards~\cite[Sec.~5.4]{rfc8446}; HTTP/2 lets the sender pad DATA and HEADERS frames and send PING frames, which the receiver's HTTP layer consumes~\cite[Secs.~6.1, 6.7]{rfc9113}. These mechanisms can shape the defender's own request bytes and timing. A padding-only record does not request an application response; PING requests a protocol acknowledgement with an eight-octet payload. Neither mechanism by itself specifies an application-shaped response or its service latency.

Cover termination therefore depends on the private signal and the provider interface. Request-side padding leaves application response activity to be addressed separately (Table~\ref{tab:cover-necessity}). A cooperating provider can synthesize padded, delayed responses without executing the business operation. Where that interface is unavailable, additional application calls are one way to generate response activity. Our experiment compares calls executed on isolated clones; it is not an impossibility result for all other defenses. A recipient that operates the service sees every exchange it receives, so against it cover must also carry plausible arguments. Reaching the application does not by itself require a business write: a read-only or no-op call elicits a real response. Cover must perform a state-changing operation only if such operations produce responses or latencies that reads cannot imitate. We call any cover that reaches the application \emph{executed cover}; the design space in Section~\ref{sec:design} separates how much business work it performs.

\begin{table}[t]\centering\footnotesize
\caption{What each terminator hides from an observer of the final hop to a third-party service. Protocol-level padding is discarded by the service's TLS or HTTP/2 layer without its cooperation~\cite{rfc8446,rfc9113}, but responses still follow real application calls.}\label{tab:cover-necessity}
\setlength{\tabcolsep}{3pt}
\begin{tabularx}{\columnwidth}{@{}P{.3\columnwidth}P{.22\columnwidth}Y@{}}\toprule
Cover mechanism & Terminated by & Still visible to the observer \\\midrule
Defender tunnel or relay & Defender endpoint & Everything on the relay's upstream hop \\\addlinespace[2pt]
TLS record padding, padding-only records & Service TLS layer & Response presence, size and latency of real calls \\\addlinespace[2pt]
HTTP/2 frame padding, PING & Service HTTP/2 layer & As above; PING acknowledges an 8-octet payload \\\addlinespace[2pt]
Read-only or no-op call & Application & Read-versus-write response differences; quota use \\\addlinespace[2pt]
Business operation (sandbox or live) & Application & Only what the attack can still separate \\\bottomrule
\end{tabularx}\end{table}


\paragraph{Observer in this paper}
Our measurements use a transit observer on the final hop to seven tool services, with no defender relay in between, and a privacy target that includes the exchange sequence and response timing. It receives destinations, ciphertext bytes and TLS record completion times. Section~\ref{sec:dispatch} additionally supplies episode and call windows, an evaluator privilege that Section~\ref{sec:noboundary} withdraws. The recipient view, where decoy arguments are also visible, is part of the threat model but not of these measurements (Section~\ref{sec:limits}). The adversary may train on traffic generated under the same cover policy; the attack suite therefore includes classifiers and cover-removal models fitted on protected training executions, alongside attacks frozen on unprotected traffic.
